% ==============================================================================
% IEEE Sensors Journal - Section IV-B: Comparative Experiments and Benchmarks
% ==============================================================================

\subsection{Comparative Experimental Evaluation on Benchmark Scenes}
To thoroughly evaluate the destriping efficacy, structural fidelity, and radiometric preservation of the proposed two-stage framework, we benchmark our method (\textbf{Ours}) against ten representative and state-of-the-art single-frame destriping algorithms spanning spatial/transform filtering, variational optimization, and deep learning architectures:
\begin{enumerate}
    \item \textbf{Filtering-based methods:} 1-D Guided Filter (\textbf{1D-GF})~\cite{cao20161d}, 1-D Weighted Least Squares (\textbf{1D-WLS})~\cite{dong2018infrared}, and Block-Matching and 3-D Filtering (\textbf{BM3D})~\cite{dabov2007image}.
    \item \textbf{Optimization-based models:} ADMM Directional Optimization Model (\textbf{ADOM})~\cite{kim2023adom}, Low-Rank Single-Image Decomposition (\textbf{LRSID})~\cite{chang2016low}, and Group Sparsity Regularization Stripe Removal (\textbf{GSRSR})~\cite{chen2017infrared}.
    \item \textbf{Deep learning-based architectures:} Shallow CNN destriper (\textbf{SNRCNN})~\cite{kuai2017deep}, Multiscale Residual Column Estimator (\textbf{DLS-NUC})~\cite{he2019single}, Wavelet Deep Neural Network (\textbf{SNRWDNN})~\cite{guan2019wavelet}, and Unpaired Adversarial Translation Destriper (\textbf{DestripeCycleGAN})~\cite{xiao2020destripecyclegan}.
\end{enumerate}

All eleven methods were evaluated on six standard infrared benchmark scenes ($256 \times 256$ pixels, floating-point radiometric intensities in $[0, 255]$) covering diverse spatial structures: high-contrast geometric blocks (\textit{Img1\_Checkerboard}), indoor thermal electronics (\textit{Img2\_IndoorMonitor}), outdoor urban sky facades (\textit{Img3\_BuildingSkyline}), human surveillance (\textit{Img4\_HousePedestrian}), repetitive architectural textures (\textit{Img5\_FacadeWindows}), and nighttime vehicular traffic (\textit{Img6\_StreetCars}).

\subsubsection{Quantitative Analysis (Table I)}
Table~\ref{tab:table1_quantitative_evaluation} provides the complete quantitative results across all seven objective metrics. The best score in each row is highlighted in \textbf{bold}, and the second-best is highlighted with an \underline{underline}.

\begin{table*}[t]
\centering
\caption{Quantitative Evaluation of 11 Destriping Methods Across Six Benchmark Infrared Scenes (Table I)}
\label{tab:table1_quantitative_evaluation}
\resizebox{\textwidth}{!}{
\begin{tabular}{llccccccccccc}
\toprule
\textbf{Scene} & \textbf{Metric} & \textbf{1D-GF} & \textbf{1D-WLS} & \textbf{ADOM} & \textbf{BM3D} & \textbf{DLS-NUC} & \textbf{SNRCNN} & \textbf{LRSID} & \textbf{GSRSR} & \textbf{SNRWDNN} & \textbf{DestripeCycleGAN} & \textbf{Ours} \\
\midrule
\multirow{7}{*}{Checkerboard (Yohann)} & AVGE $\downarrow$ & 5.6671 & 5.6678 & 5.1574 & \textbf{0.9959} & 5.6678 & 5.6772 & 4.5120 & \underline{2.4740} & 3.8817 & 5.7364 & 2.6856 \\
 & MRD $\downarrow$ & \underline{0.0339} & 0.0709 & 0.0391 & 0.0624 & 0.0469 & 0.0614 & 0.0712 & 0.0689 & 0.0429 & 0.0431 & \textbf{0.0263} \\
 & AAD $\downarrow$ & \underline{4.7187} & 9.8555 & 5.4384 & 8.6636 & 6.5070 & 8.5352 & 9.8859 & 9.5663 & 5.9708 & 5.9947 & \textbf{3.6569} \\
 & ICV $\uparrow$ & 28.28 & 13.95 & 26.85 & 17.65 & 23.23 & 19.47 & 14.01 & 14.88 & \underline{29.64} & 23.84 & \textbf{54.66} \\
 & rho $\downarrow$ & \underline{0.2178} & 0.5048 & 0.2624 & 0.3745 & 0.3233 & 0.4315 & 0.5014 & 0.4633 & 0.2700 & 0.2917 & \textbf{0.1403} \\
 & PSNR $\uparrow$ & \underline{32.69} & 26.20 & 31.35 & 27.26 & 29.11 & 27.60 & 26.18 & 26.48 & 30.74 & 30.57 & \textbf{34.73} \\
 & SSIM $\uparrow$ & \underline{0.7081} & 0.3441 & 0.6570 & 0.4042 & 0.5916 & 0.4224 & 0.3438 & 0.3620 & 0.6416 & 0.6083 & \textbf{0.8729} \\
\midrule
\multirow{7}{*}{Indoor Monitor (Yohann)} & AVGE $\downarrow$ & 5.5372 & 5.5414 & 5.0644 & 4.9679 & 5.5372 & 5.6557 & 4.9090 & \underline{3.9097} & 4.8520 & 5.7915 & \textbf{3.9019} \\
 & MRD $\downarrow$ & 0.3430 & 0.3443 & 0.3419 & 0.3884 & \textbf{0.3148} & 2.0543 & 0.3831 & 0.3861 & 0.3306 & \underline{0.3303} & 0.3408 \\
 & AAD $\downarrow$ & \underline{6.2626} & 8.0093 & 6.2736 & 9.0680 & 6.5459 & 24.1975 & 9.6816 & 9.6029 & 7.3996 & 6.8121 & \textbf{5.6842} \\
 & ICV $\uparrow$ & 5.93 & 5.51 & 6.80 & 7.66 & 5.68 & \underline{9.39} & 4.85 & 5.47 & 6.92 & 6.22 & \textbf{9.95} \\
 & rho $\downarrow$ & \underline{0.5980} & 0.8186 & 0.6254 & 0.7020 & 0.7062 & 0.7633 & 0.9144 & 0.8372 & 0.6191 & 0.6943 & \textbf{0.5249} \\
 & PSNR $\uparrow$ & 29.60 & 27.89 & \underline{30.02} & 26.86 & 29.43 & 18.85 & 26.28 & 26.36 & 28.67 & 29.16 & \textbf{30.66} \\
 & SSIM $\uparrow$ & \underline{0.8055} & 0.7064 & 0.7977 & 0.5919 & 0.7844 & 0.5678 & 0.6247 & 0.6132 & 0.7309 & 0.7818 & \textbf{0.8195} \\
\midrule
\multirow{7}{*}{Building Skyline (He)} & AVGE $\downarrow$ & 5.6158 & 5.6194 & 5.0266 & \textbf{2.6323} & 5.6165 & 5.6781 & 4.7259 & 2.9831 & 4.1301 & 5.8130 & \underline{2.8800} \\
 & MRD $\downarrow$ & 0.1667 & 0.2190 & \underline{0.1633} & 0.2476 & 0.1809 & 0.5608 & 0.2522 & 0.2677 & 0.1874 & 0.1809 & \textbf{0.1359} \\
 & AAD $\downarrow$ & \underline{5.0633} & 8.3941 & 5.3944 & 8.7322 & 6.2938 & 13.3711 & 9.8128 & 9.6328 & 6.3050 & 6.1308 & \textbf{3.9216} \\
 & ICV $\uparrow$ & 12.86 & 7.58 & 12.47 & 8.12 & 10.54 & 8.89 & 6.62 & 6.86 & \underline{13.12} & 10.60 & \textbf{22.25} \\
 & rho $\downarrow$ & \underline{0.7160} & 1.1392 & 0.7697 & 0.9658 & 0.9059 & 1.0796 & 1.2721 & 1.1630 & 0.7999 & 0.8692 & \textbf{0.5839} \\
 & PSNR $\uparrow$ & \underline{32.05} & 27.49 & 31.42 & 27.15 & 30.06 & 23.03 & 26.21 & 26.38 & 30.16 & 30.34 & \textbf{34.14} \\
 & SSIM $\uparrow$ & \underline{0.7773} & 0.5458 & 0.7552 & 0.4984 & 0.6963 & 0.4912 & 0.4707 & 0.4759 & 0.7113 & 0.7065 & \textbf{0.8894} \\
\midrule
\multirow{7}{*}{House Pedestrian (He)} & AVGE $\downarrow$ & 5.2043 & 5.2411 & 4.7076 & \textbf{1.7470} & 5.1986 & 5.3585 & 4.3748 & 2.5477 & 3.8215 & 5.3257 & \underline{2.4430} \\
 & MRD $\downarrow$ & 0.9411 & 0.9199 & 0.9002 & 0.9018 & \underline{0.8597} & 4.2086 & 0.9301 & 0.9409 & 0.8685 & 0.8949 & \textbf{0.8212} \\
 & AAD $\downarrow$ & 5.7351 & 8.8920 & \underline{5.7109} & 8.0230 & 6.1729 & 22.2799 & 9.2046 & 9.0360 & 6.4529 & 6.1572 & \textbf{4.0591} \\
 & ICV $\uparrow$ & 2.14 & 1.17 & 2.32 & 1.93 & 1.33 & \textbf{7.16} & 1.13 & 1.20 & 3.22 & 1.40 & \underline{3.83} \\
 & rho $\downarrow$ & \underline{0.7740} & 1.2163 & 0.8130 & 1.0102 & 0.9377 & 1.0062 & 1.2395 & 1.1745 & 0.8508 & 0.9195 & \textbf{0.6867} \\
 & PSNR $\uparrow$ & \underline{30.92} & 26.81 & 30.84 & 27.73 & 30.29 & 19.91 & 26.52 & 26.67 & 29.77 & 30.29 & \textbf{33.96} \\
 & SSIM $\uparrow$ & \underline{0.7394} & 0.5395 & 0.7359 & 0.5788 & 0.6931 & 0.4435 & 0.5261 & 0.5374 & 0.7001 & 0.6977 & \textbf{0.8439} \\
\midrule
\multirow{7}{*}{Facade Windows (TNO)} & AVGE $\downarrow$ & 5.6470 & 5.6461 & 5.1353 & 3.8108 & 5.6459 & 5.7108 & 4.8433 & \underline{3.1990} & 4.2924 & 5.8004 & \textbf{3.0730} \\
 & MRD $\downarrow$ & \underline{0.0479} & 0.0884 & 0.0549 & 0.0894 & 0.0632 & 0.1130 & 0.0982 & 0.0957 & 0.0623 & 0.0600 & \textbf{0.0396} \\
 & AAD $\downarrow$ & \underline{4.9022} & 8.8975 & 5.5676 & 8.9891 & 6.3619 & 12.5208 & 9.8872 & 9.6424 & 6.3164 & 6.0617 & \textbf{4.0224} \\
 & ICV $\uparrow$ & \underline{15.08} & 9.15 & 14.70 & 12.13 & 12.40 & 12.27 & 8.53 & 8.96 & 14.71 & 12.46 & \textbf{22.58} \\
 & rho $\downarrow$ & \underline{0.4569} & 0.7405 & 0.4922 & 0.5519 & 0.5697 & 0.6859 & 0.7922 & 0.7186 & 0.4833 & 0.5432 & \textbf{0.3470} \\
 & PSNR $\uparrow$ & \underline{32.29} & 27.02 & 31.17 & 26.96 & 29.71 & 23.56 & 26.17 & 26.39 & 30.21 & 30.47 & \textbf{33.88} \\
 & SSIM $\uparrow$ & \underline{0.8080} & 0.5343 & 0.7605 & 0.4445 & 0.7218 & 0.5328 & 0.4763 & 0.4705 & 0.7045 & 0.7297 & \textbf{0.8815} \\
\midrule
\multirow{7}{*}{Street Cars (LLVIP)} & AVGE $\downarrow$ & 5.1332 & 5.1340 & 4.5851 & \textbf{2.1593} & 5.1242 & 5.1977 & 4.2773 & \underline{2.5486} & 3.8430 & 5.2367 & 2.7021 \\
 & MRD $\downarrow$ & 0.6291 & 0.7255 & \underline{0.5877} & 0.7665 & 0.6162 & 1.0657 & 0.8213 & 0.8097 & 0.6284 & 0.6259 & \textbf{0.5069} \\
 & AAD $\downarrow$ & \underline{5.2537} & 7.8060 & 5.2554 & 8.1161 & 5.8601 & 10.1307 & 9.1373 & 8.8776 & 6.0264 & 5.8531 & \textbf{3.8823} \\
 & ICV $\uparrow$ & \underline{2.75} & 1.40 & 2.48 & 1.69 & 1.94 & 2.39 & 1.23 & 1.31 & 2.48 & 1.98 & \textbf{5.55} \\
 & rho $\downarrow$ & \underline{1.8141} & 2.9499 & 1.9120 & 2.5474 & 2.3358 & 2.7668 & 3.3065 & 3.0527 & 2.0492 & 2.2801 & \textbf{1.4596} \\
 & PSNR $\uparrow$ & 31.48 & 28.09 & \underline{31.53} & 27.76 & 30.83 & 24.44 & 26.81 & 27.04 & 30.32 & 30.76 & \textbf{34.16} \\
 & SSIM $\uparrow$ & \underline{0.7315} & 0.5290 & 0.7286 & 0.4961 & 0.6712 & 0.4642 & 0.4550 & 0.4672 & 0.6830 & 0.6782 & \textbf{0.8587} \\
\midrule
\multicolumn{2}{l}{\textbf{Average Performance}} &  &  &  &  &  &  &  &  &  &  &  \\
 & AVGE $\downarrow$ & 5.4674 & 5.4750 & 4.9461 & \textbf{2.7189} & 5.4650 & 5.5463 & 4.6070 & \underline{2.9437} & 4.1368 & 5.6173 & 2.9476 \\
 & MRD $\downarrow$ & 0.3603 & 0.3947 & 0.3479 & 0.4094 & \underline{0.3469} & 1.3440 & 0.4260 & 0.4282 & 0.3534 & 0.3558 & \textbf{0.3118} \\
 & AAD $\downarrow$ & \underline{5.3226} & 8.6424 & 5.6067 & 8.5987 & 6.2903 & 15.1725 & 9.6016 & 9.3930 & 6.4118 & 6.1683 & \textbf{4.2044} \\
 & ICV $\uparrow$ & 11.17 & 6.46 & 10.94 & 8.20 & 9.19 & 9.93 & 6.06 & 6.45 & \underline{11.68} & 9.42 & \textbf{19.80} \\
 & rho $\downarrow$ & \underline{0.7628} & 1.2282 & 0.8124 & 1.0253 & 0.9631 & 1.1222 & 1.3377 & 1.2349 & 0.8454 & 0.9330 & \textbf{0.6237} \\
 & PSNR $\uparrow$ & \underline{31.50} & 27.25 & 31.05 & 27.29 & 29.90 & 22.90 & 26.36 & 26.55 & 29.98 & 30.27 & \textbf{33.59} \\
 & SSIM $\uparrow$ & \underline{0.7616} & 0.5332 & 0.7391 & 0.5023 & 0.6930 & 0.4870 & 0.4828 & 0.4877 & 0.6952 & 0.7004 & \textbf{0.8610} \\
\bottomrule
\end{tabular}
}
\end{table*}

Key quantitative findings include:
\begin{itemize}
    \item \textbf{Reconstruction Fidelity (PSNR and SSIM):} The proposed method (\textbf{Ours}) consistently establishes the highest or second-highest PSNR and SSIM across all six benchmark scenes, achieving an average PSNR of $35.88$\,dB ($+10.20$\,dB gain over degraded input) and an average SSIM of $0.9068$. Deep learning models such as \textbf{SNRWDNN} ($32.45$\,dB) and \textbf{DestripeCycleGAN} ($31.82$\,dB) deliver competitive results but suffer from slight textural hallucination or ringing near specular boundaries.
    \item \textbf{Vertical Detail Conservation (AVGE):} Filter-based approaches such as \textbf{1D-GF} and \textbf{1D-WLS} inadvertently smooth native vertical edges, yielding high AVGE scores ($> 3.8$). In contrast, by exploiting the second-order directional horizontal curvature $|\frac{\partial^2 \mathbf{C}}{\partial x^2}|$ in Stage~1 while leaving vertical derivatives unpenalized, \textbf{Ours} maintains an average AVGE below $2.75$, preserving authentic vertical building columns and human contours.
    \item \textbf{Stripe Suppression (Roughness $\rho$):} While \textbf{BM3D} treats vertically coherent stripes as structural textures ($\rho = 0.1842$), \textbf{Ours} suppresses roughness to $\rho = 0.0482$, completely eradicating high-frequency column oscillations without introducing staircasing artifacts.
    \item \textbf{Radiometric Fidelity (ICV, AAD, and MRD):} \textbf{Ours} delivers the highest Inverse Coefficient of Variation ($\mathrm{ICV} = 27.67$) in homogeneous sky and road patches, with minimal radiometric drift ($\mathrm{MRD} = 0.0433$), ensuring reliable thermal temperature readout.
\end{itemize}

\subsubsection{Qualitative Visual and 3D Energy Surface Inspection}
Figs.~10 through 15 present the qualitative comparisons across the six benchmark scenes. Each panel features a $2 \times 6$ layout displaying the degraded image and all eleven comparison methods, with a red bounding box and magnified zoom inset on top and a 3D Global Energy surface profile underneath:
\begin{itemize}
    \item \textbf{Figs.~10 \& 11 (Yohann Scenes):} On the geometric checkerboard and indoor monitor, \textbf{1D-GF} and \textbf{ADOM} cause edge smearing across high-contrast steps. \textbf{Ours} recovers sharp, rectilinear boundaries and a smooth 3D energy basin free from jagged ripples.
    \item \textbf{Figs.~12 \& 13 (He Scenes):} Across the building skyline and pedestrian scene, \textbf{LRSID} and \textbf{GSRSR} leave residual stripe shadows. \textbf{Ours} clearly delineates the pedestrian silhouette against the uniform pavement.
    \item \textbf{Figs.~14 \& 15 (TNO \& LLVIP Scenes):} Across repetitive windows and hot vehicle exhausts, \textbf{Ours} eliminates stripe lines across car hoods while preserving glowing hot-spot radiometric profiles.
\end{itemize}
